% ============================================================
%                     INPUT AND LANGUAGE
% ============================================================

% Modern LaTeX already uses UTF-8 by default.
\usepackage[utf8]{inputenc} %Il pacchetto traduce varie codifiche di input standard e di altro tipo in un "linguaggio interno LATEX".

% ============================================================
%                         GRAPHICS
% ============================================================

\usepackage{graphicx}%potenziamento del pacchetto graphics con opzioni aggiuntive al commando \includegraphics
\usepackage{animate}

\graphicspath{
	{"D:/Ale/Universita/Triennale (Bachelor)/Libri università/LaTeX/progetti/file base"}
	{"D:/Ale/Universita/Triennale (Bachelor)/Libri università/LaTeX/progetti/file base/tex/latex"}
	{"D:/Ale/Universita/Master/Thesis/3) Physics 30-70 Computation - Machine Learning documents/Neural Networks/Networks 	and Deep Learning"}
	{../Images/}
    {../Images/Chapter 1/}
    {../Images/Chapter 2/}
    {../Images/Chapter 3/}
	%{"D:/Ale/Universita/Master/Thesis/3) Physics 30-70 Computation - Machine Learning documents/Neural Networks/Networks 	and Deep Learning/Images/Chapter 1"}
	%{"D:/Ale/Universita/Master/Thesis/3) Physics 30-70 Computation - Machine Learning documents/Neural Networks/Networks 	and Deep Learning/Images/Chapter 2"}
}

% ============================================================
%                       LANGUAGE
% ============================================================

\usepackage[english]{babel} %cambia i termini viaulizzati nella lingua scelta esempio references invece che bibliografia.


% ============================================================
%                    MATHEMATICS
% ============================================================

\usepackage{amsmath} %pacchetto per utilizzo di equazioni matematiche
\usepackage{amsfonts} %aggiunta di simboli matematici
\usepackage{amssymb} %aggiunge ulteriori simboli matematici.
\usepackage{mathtools} %fornisce una serie di pacchetti progettati per migliorare l'aspetto dei documenti contenenti molta matematica.
\usepackage{mathrsfs} %introduce un nuovo tipo di font di tipo matematico
\usepackage{nccmath} %aggiunge alcuni commandi matematici
\usepackage{bigints} %permetter di scrivere integrali piu grandi
\usepackage{cancel}        %il pacchetto introduce simboli di semplificazione delle equazioni \bacancel{}
\usepackage{old-arrows} % Questo pacchetto rende disponibili le frecce Computer Modern vecchio stile.
\usepackage{tcolorbox}
% ============================================================
%                         FIGURES
% ============================================================

\usepackage{caption}       %permette di customizzare le didascalie
\usepackage{float} %migliora l'interfaccia per descrivere oggetti fluttuanti come figure e tabelle.
\usepackage{afterpage} %serve a itrodurre il comando afterpage per sistemare le pagine con figure e testi
\usepackage[rightcaption]{sidecap}  %il pacchetto permette di posizionare la didascalia a destra dell'immagine
\usepackage{subfigure} %Fornisce supporto per la manipolazione e il riferimento di figure e tabelle piccole o "sottostanti" all'interno di una singola figura o ambiente di tabella
\usepackage{wrapfig} % Consente alle figure o alle tabelle di avere il testo avvolto attorno ad esse

% ============================================================
%                         TABLES
% ============================================================

\usepackage{makecell}  %permette di modificare le tabelle
\usepackage{tabularx} %permette di definire un ambiente tabularx, una estensione di tabular, che ha un aggiuntivo designatore di colonne.

% ============================================================
%                    CIRCUITS AND PLOTS
% ============================================================

\usepackage{circuitikz}  %il pacchetto permette di disegnare circuiti
\usepackage{pgfplots} %disegna grafci ad alta qualità di funzioni in scala normale o logaritmica
\usepackage{pgfplotstable}% is used to read, process, format, and typeset tabular data—especially data stored in .csv or other text files.
\usepackage{tkz-fct} %offre facile accesso ai grafici di programmazione delle funzioni con TikZ e gnuplot.

% ============================================================
%                    DOCUMENT STRUCTURE
% ============================================================

\usepackage[toc,page]{appendix} %aggiunge alla tabella dei contenuti l'appendice e provvede alcuni commandi per l'appendice
\usepackage[nottoc]{tocbibind} %permette di aggiungere la bibliografia alla tabella dei contenuti
\usepackage{comment}  %permette di commentare intere sezioni invece che frasi singole
\usepackage{enumerate} %permette di modificare le liste di punti numerati


% ============================================================
%                       PAGE LAYOUT
% ============================================================

\usepackage[margin=30mm]{geometry}  %permette di impostare i margini del foglio otre che il tipo di font e altri dettagli di impaginazione
\usepackage{fancyhdr} %fornisce funzionalità estese, sia per la creazione di intestazioni e piè di pagina, sia per il controllo del loro utilizzo

% ============================================================
%                    FONTS AND DECORATION
% ============================================================

\usepackage{calligra}      	%pacchetto che introduce i font calligra
\usepackage{miama}         	%pacchetto che introduce i font miama

\usepackage{lmodern}			%cambia il font di tutto il docuemnto

%\usepackage{stix2}

%\usepackage{tgtermes}


% ============================================================
%                    CODE LISTINGS
% ============================================================

\usepackage{listings}

% ============================================================
%                      OTHER PACKAGES
% ============================================================

\usepackage{authblk} %Il pacchetto ridefinisce il comando \author in modo che funzioni normalmente o per consentire uno stile di nota a piè di pagina.
\usepackage{afterpage}
\usepackage{lipsum}  %genera dummy text

% ============================================================
%                     HYPERLINKS
% ============================================================
%\usepackage[hidelinks]{hyperref}
\usepackage[hidelinks, pdfstartview=FitH]{hyperref}
%\usepackage{hyperref} %permette di creare link ipertesto (il comando [hidelinks] nasconde i controni rossi


% ============================================================
%                         SETTINGS
% ============================================================

% ---------- Hyperref ----------

\hypersetup{
	colorlinks=true,
	linkcolor={red!50!black},
	citecolor={blue!50!black},
	urlcolor={blue!80!black}
}


% ---------- Code listings ----------

\lstset{
    basicstyle=\ttfamily\footnotesize,
    backgroundcolor=\color{gray!10},
    frame=none,
    breaklines=true,
    escapeinside={(*@}{@*)}
}


% ---------- Equations ----------

\numberwithin{equation}{section} %permette di avere le equazioni numerati in base alla sezione o sottosezione in cui esse sono disposte


% ---------- PGFPlots ----------

\pgfplotsset{
	compat=1.16,
    width=\textwidth,
    height=0.625\textwidth
}




% ============================================================
%                     CUSTOM COMMANDS
% ============================================================

\newcommand{\quotes}[1]{``#1''}

\DeclareRobustCommand{\rchi}{{\mathpalette\irchi\relax}}   %aggiunge un comando \rhi per disgnare meglio il chi greco

\newcommand{\irchi}[2]{\raisebox{\depth}{$#1\chi$}} % inner command, used by \rchi %modi diversi per scrivere il chi greco

\newcommand{\eqstackrel}[1]{
	\stackrel{
		\substack{
			\mathclap{#1}\\[0.5ex]
			\displaystyle\uparrow\\
			 ~
		}
	}{ = }
} %permette di introdurre il commando \eqstackrel{} in maniera da disegnare un uguale con una freccia legata al testa tra parentesi

\newcommand\Item[1][]{%
  \ifx\relax#1\relax
  	\item
  \else
  	\item[#1]
  \fi
  \abovedisplayskip=0pt
  \abovedisplayshortskip=0pt
  \belowdisplayshortskip=0pt
  ~\vspace*{-\baselineskip}
}

\newcommand{\heart}{\ensuremath\heartsuit} %aggiunge un commando \heart in maniera da aggiungere un simbolo di un cuoricino

\newcommand{\verteq}{\rotatebox{90}{$\,=$}} %aggiunge un commando \verteq in mnaiera da aggiungere un segno di uguale in verticale anziche in orizzontale

\newcommand{\equalto}[2]
	{\underset{
		\scriptstyle\overset{\mkern4mu\verteq}{#2}
		}{#1}
} %aggiunge un commando \equalto{}{} in maniera da associare al testo dentro le prime graffe il testo dentro le seconde tramite un segno di uguale posizionato in verticale
%\newcommand{\superv}{Prof. William Joseph Weber} %aggiunge un ulteriore autore




%create a circuit representina Neuron in a NN
\newcommand{\neuron}[4]{%
\begin{tikzpicture}[
    >=stealth,
    thick,
    every node/.style={font=\large}
]

% Input
\node (input) at (-1,0) {$x=1$};

% Neuron
\draw (4,0) circle (1);
\node at (4,0) {$\sigma$};

% ------------------------------------------------
% Weight
% ------------------------------------------------

% First arrow
\draw[->] (-0.3,0) -- (3,0);

% Weight label and value
\node at (1.25,0.85) {$w:~#1$};

% Weight box
\fill[#4!40] (1.2,0.25) rectangle (1.2+#1*0.5,0.60);
\draw (0.2,0.25) rectangle (2.2,0.60);

% Zero line
\draw (1.2,0.25) -- (1.2,0.60);

% ------------------------------------------------
% Bias
% ------------------------------------------------


% Bias label and value
\node at (4.05,1.75) {$b:~#2$};

% Bias box
\fill[#4!40] (4.0,1.15) rectangle (4.0+#2*0.5,1.50);
\draw (3.0,1.15) rectangle (5.0,1.5);

% Zero line
\draw (4.0,1.15) -- (4.0,1.50);

% ------------------------------------------------
% Output
% ------------------------------------------------

% Output arrow
\draw[->] (5,0) -- (8,0);

% Output label and value
\node at (9.15,0.45) {$a:~#3$};

% Output box
\fill[#4!40] (8.20,-0.15) rectangle (8.2+#3*2,0.2);
\draw (8.20,-0.15) rectangle (10.20,0.2);


\end{tikzpicture}%
}


\newcommand{\Costplot}[4]{%
\pgfmathtruncatemacro{\epoch}{#2+1}

\ifnum\epoch<10
    \edef\formattedepoch{00\epoch}
\else
	\ifnum\epoch<100
    		\edef\formattedepoch{0\epoch}
	\else
    		\edef\formattedepoch{\epoch}
	\fi
\fi

\begin{tikzpicture}%[show background rectangle]
\begin{axis}[
    xlabel={$Epochs$},
    ylabel={$Cost$},
    xmin=0,
    xmax=330,
    ymin=0,
    ymax=1.1*#4,
    grid=major,
    width=7cm,
    height=4cm,
    scale only axis,
    legend style={
    		text width=2cm,
    		align=right,
    	},
]
\addplot[
    #3,
    thick,
    smooth,
] coordinates {
    #1
};
\addlegendentry{Epoch: \formattedepoch}
\end{axis}
\end{tikzpicture}%
}



\newcommand{\animateNeuron}[2]{%

    % Number of rows
    \pgfplotstablegetrowsof{#1}
    \pgfmathtruncatemacro{\lastrow}{\pgfplotsretval}

    % Animation
    \begin{animateinline}[
        controls,
    ]{60}
	

		% ------------------------------------------------
		% Loop through CSV
		% ------------------------------------------------
		\def\mydata{}
		
        \multiframe{\lastrow}{i=0+1}{

	    		% Get weight
            \pgfplotstablegetelem{\i}{w}\of{#1}
            \edef\currentw{\pgfplotsretval}

 	    		% Get bias
            \pgfplotstablegetelem{\i}{b}\of{#1}
            \edef\currentb{\pgfplotsretval}

	    		% Get output
            \pgfplotstablegetelem{\i}{a}\of{#1}
            \edef\currenta{\pgfplotsretval}
            
            % Get Cost
            \pgfplotstablegetelem{\i}{C}\of{#1}
            \edef\currentC{\pgfplotsretval}
            
            % Get inital Cost (to set yrange)
            \ifnum\i=0
            	\xdef\initialC{\currentC}
            	\fi
			
			\xdef\mydata{\mydata (\i,\currentC)}
			
			\hspace{-1cm}%
			\makebox[\linewidth][c]{%
    				\parbox{\linewidth}{%
        				\centering
        				\neuron{\currentw}{\currentb}{\currenta}{#2}

        				\vspace{0.5cm}
					
					\hspace{-2cm}%
        				\Costplot{\mydata}{\i}{#2}{\initialC}%
    				}%
			}
	    		
	    		% New animation frame
            %\ifnum\i<\lastrow
            %    \newframe
            %\fi
        }

    \end{animateinline}
}

\newcommand{\softmaxBars}[8]{%
\begin{tikzpicture}[
    >=stealth,
    thick,
    every node/.style={font=\large}
]

% z1 label and value
\node at (-2.3,2.6) {$z_1^L:~#1$};

% z1 box
\fill[red!40] (-2,3) rectangle (-2+#1*0.2,3.5);
\draw (-3,3) rectangle (-1,3.5);

% Zero line
\draw (-2,3) -- (-2,3.5);



% z2 label and value
\node at (-2.3,.6) {$z_2^L:~#2$};

% z2 box
\fill[red!40] (-2,1) rectangle (-2+#2*0.2,1.5);
\draw (-3,1) rectangle (-1,1.5);

% Zero line
\draw (-2,1) -- (-2,1.5);



% z3 label and value
\node at (-2.3,-1.4) {$z_3^L:~#3$};

% z3 box
\fill[red!40] (-2,-0.5) rectangle (-2+#3*0.2,-1);
\draw (-3,-0.5) rectangle (-1,-1);

% Zero line
\draw (-2,-0.5) -- (-2,-1);



% z4 label and value
\node at (-2.3,-3.4) {$z_4^L:~#4$};

% z4 box
\fill[red!40] (-2,-2.5) rectangle (-2+#4*0.2,-3);
\draw (-3,-2.5) rectangle (-1,-3);

% Zero line
\draw (-2,-2.5) -- (-2,-3);






% a1 label and value
\node at (1.7,2.6) {$a_1^L:~#5$};

% a1 box
\fill[blue!40] (1,3) rectangle (1+#5*2,3.5);
\draw (1,3) rectangle (3,3.5);



% a2 label and value
\node at (1.7,.6) {$a_2^L:~#6$};

% a2 box
\fill[blue!40] (1,1) rectangle (1+#6*2,1.5);
\draw (1,1) rectangle (3,1.5);



% a3 label and value
\node at (1.7,-1.4) {$a_3^L:~#7$};

% a3 box
\fill[blue!40] (1,-0.5) rectangle (1+#7*2,-1);
\draw (1,-0.5) rectangle (3,-1);



% a4 label and value
\node at (1.7,-3.4) {$a_4^L:~#8$};

% a4 box
\fill[blue!40] (1,-2.5) rectangle (1+#8*2,-3);
\draw (1,-2.5) rectangle (3,-3);


\end{tikzpicture}%
}

\newcommand{\animateBars}[1]{%

    % Number of rows
    \pgfplotstablegetrowsof{#1}
    \pgfmathtruncatemacro{\lastrow}{\pgfplotsretval}

    % Animation
    \begin{animateinline}[
        controls,
    ]{30}
	

		% ------------------------------------------------
		% Loop through CSV
		% ------------------------------------------------
				
        \multiframe{\lastrow}{i=0+1}{

	    		% Get z1
            \pgfplotstablegetelem{\i}{z1}\of{#1}
            \edef\currentzone{\pgfplotsretval}

 	    		% Get z2
            \pgfplotstablegetelem{\i}{z2}\of{#1}
            \edef\currentztwo{\pgfplotsretval}
            
            % Get z3
            \pgfplotstablegetelem{\i}{z3}\of{#1}
            \edef\currentzthree{\pgfplotsretval}
            
            % Get z4
            \pgfplotstablegetelem{\i}{z4}\of{#1}
            \edef\currentzfour{\pgfplotsretval}
			
			% Get a1
            \pgfplotstablegetelem{\i}{a1}\of{#1}
            \edef\currentaone{\pgfplotsretval}
            
            % Get a2
            \pgfplotstablegetelem{\i}{a2}\of{#1}
            \edef\currentatwo{\pgfplotsretval}
            
            % Get a3
            \pgfplotstablegetelem{\i}{a3}\of{#1}
            \edef\currentathree{\pgfplotsretval}
            
            % Get a4
            \pgfplotstablegetelem{\i}{a4}\of{#1}
            \edef\currentafour{\pgfplotsretval}
            
			
			\hspace{-1cm}%
			\makebox[\linewidth][c]{%
    				\parbox{\linewidth}{%
        				\centering
        				\softmaxBars{\currentzone}{\currentztwo}{\currentzthree}{\currentzfour}{\currentaone}{\currentatwo}{\currentathree}{\currentafour}

    				}%
			}
	    		
	    		% New animation frame
            %\ifnum\i<\lastrow
            %    \newframe
            %\fi
        }

    \end{animateinline}
}


% ============================================================
%                       WIDEBAR COMMAND
% ============================================================


\makeatletter %tutto il pezzo fino a \makeatother

\newcommand*\rel@kern[1]{\kern#1\dimexpr\macc@kerna}

\newcommand*\widebar[1]{%
  \begingroup
  \def\mathaccent##1##2{%
    \rel@kern{0.8}%
    \overline{\rel@kern{-0.8}\macc@nucleus\rel@kern{0.2}}%
    \rel@kern{-0.2}%
  }%
  \macc@depth\@ne
  \let\math@bgroup\@empty \let\math@egroup\macc@set@skewchar
  \mathsurround\z@ \frozen@everymath{\mathgroup\macc@group\relax}%
  \macc@set@skewchar\relax
  \let\mathaccentV\macc@nested@a
  \macc@nested@a\relax111{#1}%
  \endgroup
}
\makeatother %tutto il pezzo da \makeatletter fino a \makeatother introduce il comando \widebar{} per avere una barra piu larga sopra il testo scritto tra le graffe

\newcommand{\gigi}[1]{#1}

% ============================================================
%                       TITLE PAGE
% ============================================================


\makeatletter

% Default: no cover image
\def\@coverimage{}

% User command
\newcommand{\coverimage}[2]{%
    \def\@coverimage{%
        \includegraphics[width=#2\linewidth]{#1}%
    }%
}
\makeatother

\makeatletter

% Default: no cover image
\def\@citeimage{}

% User command
\newcommand{\citeimage}[1]{%
    \def\@citeimage{{\footnotesize #1}}%
}
\makeatother


\makeatletter

\renewcommand\maketitle{
\begin{titlepage}
    \centering
    \scshape

    \vspace*{1.5\baselineskip}

    % ===================
    %   Title Section
    % ===================
	
	% Double line above the title
    \rule{\linewidth}{1.6pt}\\[-2pt]
    \rule{\linewidth}{0.4pt}

    \vspace{0.75\baselineskip}
	
	% Title
    {\fontsize{26pt}{31pt}\selectfont \@title\par}

    \vspace{0.75\baselineskip}
	
	% Double line below the title
    \rule{\linewidth}{0.4pt}\\[-2pt]
    \rule{\linewidth}{1.6pt}
	
	\vspace{2.5\baselineskip}
    %\vspace{5\baselineskip}

    % =================
    %   Cover Image
    % =================

    \@coverimage

    \begin{flushleft}
        {\fontsize{9pt}{11pt}\selectfont \@citeimage}
    \end{flushleft}

    \vfill

    % =================
    %   Author
    % =================

    {\fontsize{23pt}{28pt}\selectfont \@author\par}

    \vspace{1.5\baselineskip}

\end{titlepage}
}

\makeatother

% ============================================================
%                       PAGE STYLE
% ============================================================

% per impostare l'intestazione e il piè di pagina prima si set lo style a fancy con \pagestyle{fancy}, dopo si pulisce l'intestazione e il piè di pagina con \fancyhf{} altrimenti in seguito si avrebbe i testi di deafult che continuerebbero ad apparire, e infine gli ultimi tre comandi aggiungono il numero della pagina (a destra per i numeri dispari e a sinistra per quelli pari), il testo della sezione in alto (a sinistra per pagine pari) e il testo del paragrafo (a sinistra per pagine dispari).

\pagestyle{fancy}

\fancyhf{}

\fancyfoot[RO,LE]{\thepage}

\fancyhead[LE]{\nouppercase{\leftmark}}

\fancyhead[LO]{\rightmark}







% ============================================================
%                  OPTIONAL / UNUSED PACKAGES
% ============================================================



% Italian language (alternative to English)
%\usepackage[italian]{babel}

% Better looking tables
%\usepackage{booktabs}     %il pacchetto migliora le tabelle

% Advanced float/caption positioning
% \usepackage[
% 	  capbesideposition={top,outside},
% 	  facing=yes,
%	  capbesidesep=quad
% ]{floatrow}

% Enhanced underlining
%\usepackage{soulutf8} %permette di migliorare le funzioni di \underline usando al posto suo il comando \ul

% Table of contents customization
%\usepackage{titletoc}
%\usepackage{tocloft}